\documentclass[11pt,a4paper]{article}

% Paquetes básicos
\usepackage[utf8]{inputenc}
% Spanish babel removed for compatibility
\usepackage{amsmath,amssymb,amsthm}
\usepackage{graphicx}
\usepackage{float}
\usepackage{booktabs}
\usepackage{geometry}
\geometry{margin=2.5cm}
\usepackage{hyperref}
\hypersetup{
    colorlinks=true,
    linkcolor=blue,
    urlcolor=blue,
    citecolor=blue
}

% Paquetes para código
\usepackage{listings}
\usepackage{xcolor}

% Configuración de listings para Python
\lstdefinestyle{python}{
    language=Python,
    basicstyle=\ttfamily\small,
    keywordstyle=\color{blue}\bfseries,
    commentstyle=\color{gray}\itshape,
    stringstyle=\color{red},
    showstringspaces=false,
    numbers=left,
    numberstyle=\tiny\color{gray},
    stepnumber=1,
    numbersep=8pt,
    frame=single,
    breaklines=true,
    breakatwhitespace=true,
    tabsize=4,
    captionpos=b
}

% Configuración de listings para R
\lstdefinestyle{rstyle}{
    language=R,
    basicstyle=\ttfamily\small,
    keywordstyle=\color{blue}\bfseries,
    commentstyle=\color{gray}\itshape,
    stringstyle=\color{red},
    showstringspaces=false,
    numbers=left,
    numberstyle=\tiny\color{gray},
    stepnumber=1,
    numbersep=8pt,
    frame=single,
    breaklines=true,
    breakatwhitespace=true,
    tabsize=4,
    captionpos=b
}

% Teoremas y definiciones
\theoremstyle{definition}
\newtheorem{definicion}{Definición}[section]
\newtheorem{ejemplo}{Ejemplo}[section]
\newtheorem{teorema}{Teorema}[section]

% Información del documento
\title{\textbf{Unidad 4: ANOVA y Diseño de Experimentos} \\ 
\Large Métodos de Análisis de Datos 1}
\author{Departamento de Matemática \\ Universidad Nacional del Sur}
\date{Primer Semestre 2026}

\begin{document}

\maketitle

\tableofcontents
\newpage

\section{Introducción}

El Análisis de Varianza (ANOVA, por sus siglas en inglés \textit{Analysis of Variance}) es una técnica estadística fundamental para comparar las medias de tres o más grupos. Mientras que las pruebas t permiten comparar dos grupos, ANOVA extiende este concepto a múltiples grupos simultáneamente, controlando el error tipo I que se acumularía al realizar múltiples comparaciones por pares.

El principio básico de ANOVA es descomponer la variabilidad total de los datos en componentes: variabilidad entre grupos (explicada por el factor) y variabilidad dentro de los grupos (no explicada, o error). Si la variabilidad entre grupos es significativamente mayor que la variabilidad dentro de los grupos, hay evidencia de que al menos una media grupal difiere de las demás.

\subsection{Aplicaciones}

ANOVA se utiliza en una amplia variedad de contextos:

\begin{itemize}
    \item \textbf{Medicina:} Comparar la efectividad de diferentes tratamientos
    \item \textbf{Agronomía:} Evaluar el efecto de distintos fertilizantes en el rendimiento de cultivos
    \item \textbf{Psicología:} Analizar diferencias en desempeño cognitivo entre grupos de edad
    \item \textbf{Manufactura:} Estudiar el efecto de diferentes máquinas o procesos en la calidad del producto
    \item \textbf{Marketing:} Comparar la efectividad de diferentes estrategias publicitarias
\end{itemize}

\section{ANOVA de un Factor}

\subsection{Modelo Estadístico}

El modelo ANOVA de un factor (o de una vía) considera una variable respuesta cuantitativa $Y$ y un factor categórico con $k$ niveles o grupos. El modelo estadístico es:

\begin{equation}
Y_{ij} = \mu + \tau_i + \varepsilon_{ij}
\end{equation}

donde:
\begin{itemize}
    \item $Y_{ij}$ es la $j$-ésima observación del grupo $i$
    \item $\mu$ es la media global
    \item $\tau_i$ es el efecto del tratamiento $i$ (con $i = 1, 2, \ldots, k$)
    \item $\varepsilon_{ij} \sim N(0, \sigma^2)$ es el error aleatorio
    \item $j = 1, 2, \ldots, n_i$ (donde $n_i$ es el número de observaciones en el grupo $i$)
\end{itemize}

El modelo puede expresarse equivalentemente como:
\begin{equation}
Y_{ij} = \mu_i + \varepsilon_{ij}
\end{equation}

donde $\mu_i = \mu + \tau_i$ es la media del grupo $i$.

\subsection{Hipótesis a Contrastar}

Las hipótesis del ANOVA de un factor son:

\begin{align*}
H_0 &: \mu_1 = \mu_2 = \cdots = \mu_k \quad \text{(todas las medias son iguales)} \\
H_1 &: \text{Al menos una media difiere de las demás}
\end{align*}

Equivalentemente, en términos de los efectos:
\begin{align*}
H_0 &: \tau_1 = \tau_2 = \cdots = \tau_k = 0 \\
H_1 &: \text{Al menos un } \tau_i \neq 0
\end{align*}

\subsection{Descomposición de la Variabilidad}

ANOVA descompone la variabilidad total (SST, \textit{Sum of Squares Total}) en dos componentes:

\begin{equation}
\text{SST} = \text{SSB} + \text{SSW}
\end{equation}

donde:
\begin{itemize}
    \item \textbf{SST (Suma de Cuadrados Total):} Variabilidad total de los datos
    \begin{equation}
    \text{SST} = \sum_{i=1}^{k} \sum_{j=1}^{n_i} (Y_{ij} - \bar{Y})^2
    \end{equation}
    
    \item \textbf{SSB (Suma de Cuadrados Entre Grupos):} Variabilidad explicada por el factor
    \begin{equation}
    \text{SSB} = \sum_{i=1}^{k} n_i (\bar{Y}_i - \bar{Y})^2
    \end{equation}
    
    \item \textbf{SSW (Suma de Cuadrados Dentro de Grupos):} Variabilidad no explicada (error)
    \begin{equation}
    \text{SSW} = \sum_{i=1}^{k} \sum_{j=1}^{n_i} (Y_{ij} - \bar{Y}_i)^2
    \end{equation}
\end{itemize}

Aquí, $\bar{Y}$ es la media global y $\bar{Y}_i$ es la media del grupo $i$.

\subsection{Tabla ANOVA}

Los resultados se resumen en la tabla ANOVA:

\begin{table}[H]
\centering
\begin{tabular}{lcccc}
\toprule
Fuente & Suma de Cuadrados & GL & Cuadrados Medios & Estadístico F \\
\midrule
Entre grupos & SSB & $k-1$ & MSB = SSB/(k-1) & $F = \frac{\text{MSB}}{\text{MSW}}$ \\
Dentro de grupos & SSW & $N-k$ & MSW = SSW/(N-k) & \\
Total & SST & $N-1$ & & \\
\bottomrule
\end{tabular}
\caption{Tabla ANOVA de un factor}
\end{table}

donde $N = \sum_{i=1}^{k} n_i$ es el número total de observaciones.

\subsection{Estadístico F y Decisión}

El estadístico F compara la variabilidad entre grupos con la variabilidad dentro de grupos:

\begin{equation}
F = \frac{\text{MSB}}{\text{MSW}}
\end{equation}

Bajo $H_0$, este estadístico sigue una distribución $F_{k-1, N-k}$. Se rechaza $H_0$ si:

\begin{equation}
F > F_{k-1, N-k, \alpha}
\end{equation}

o equivalentemente, si el p-valor $< \alpha$.

\subsection{Supuestos del ANOVA}

Para que los resultados del ANOVA sean válidos, deben cumplirse tres supuestos:

\begin{enumerate}
    \item \textbf{Independencia:} Las observaciones deben ser independientes entre sí
    \item \textbf{Normalidad:} Los errores $\varepsilon_{ij}$ siguen una distribución normal (o equivalentemente, $Y_{ij}$ es normal en cada grupo)
    \item \textbf{Homocedasticidad:} La varianza es constante en todos los grupos: $\text{Var}(Y_{ij}) = \sigma^2$ para todo $i$
\end{enumerate}

\textbf{Verificación de supuestos:}
\begin{itemize}
    \item \textbf{Normalidad:} Gráfico Q-Q de residuos, test de Shapiro-Wilk
    \item \textbf{Homocedasticidad:} Test de Levene, test de Bartlett, gráfico de residuos vs. valores ajustados
\end{itemize}

\subsection{Tamaño del Efecto: $\eta^2$ (Eta Cuadrado)}

El tamaño del efecto mide la proporción de variabilidad explicada por el factor:

\begin{equation}
\eta^2 = \frac{\text{SSB}}{\text{SST}}
\end{equation}

Interpretación (reglas generales de Cohen):
\begin{itemize}
    \item $\eta^2 < 0.06$: efecto pequeño
    \item $0.06 \leq \eta^2 < 0.14$: efecto mediano
    \item $\eta^2 \geq 0.14$: efecto grande
\end{itemize}

\section{Comparaciones Múltiples Post-Hoc}

Si el ANOVA rechaza $H_0$, sabemos que al menos una media difiere, pero no sabemos cuáles. Las pruebas post-hoc permiten identificar qué pares de grupos son significativamente diferentes, controlando el error tipo I acumulado.

\subsection{Test de Tukey HSD}

El test HSD (\textit{Honestly Significant Difference}) de Tukey es el más utilizado cuando todas las comparaciones por pares son de interés y los tamaños muestrales son iguales o similares.

Para cada par de grupos $i$ y $j$, se calcula:

\begin{equation}
\text{HSD} = q_{\alpha, k, N-k} \sqrt{\frac{\text{MSW}}{n}}
\end{equation}

donde $q_{\alpha, k, N-k}$ es el valor crítico de la distribución del rango estudentizado. Dos grupos difieren significativamente si:

\begin{equation}
|\bar{Y}_i - \bar{Y}_j| > \text{HSD}
\end{equation}

\subsection{Corrección de Bonferroni}

El método de Bonferroni es más conservador y controla el error tipo I familiar ajustando el nivel de significancia:

\begin{equation}
\alpha_{\text{ajustado}} = \frac{\alpha}{m}
\end{equation}

donde $m$ es el número de comparaciones. Para $k$ grupos, $m = \binom{k}{2} = \frac{k(k-1)}{2}$.

\textbf{Ventaja:} Simple y muy conservador. \\
\textbf{Desventaja:} Puede ser demasiado conservador (baja potencia) cuando hay muchas comparaciones.

\subsection{Test de Scheffé}

El método de Scheffé es el más conservador y permite realizar cualquier tipo de contraste lineal, no solo comparaciones por pares. Es útil cuando se desean hacer comparaciones complejas o no planificadas.

El valor crítico es:

\begin{equation}
S = \sqrt{(k-1) F_{k-1, N-k, \alpha}}
\end{equation}

Dos grupos difieren si:

\begin{equation}
|\bar{Y}_i - \bar{Y}_j| > S \sqrt{\text{MSW} \left(\frac{1}{n_i} + \frac{1}{n_j}\right)}
\end{equation}

\subsection{Comparación de Métodos}

\begin{table}[H]
\centering
\small
\begin{tabular}{lp{5cm}p{5cm}}
\toprule
Método & Ventajas & Cuándo usar \\
\midrule
Tukey HSD & Potente para comparaciones por pares, controla error familiar & Comparar todos los pares, tamaños similares \\
Bonferroni & Simple, muy conservador & Pocas comparaciones planificadas \\
Scheffé & Más flexible, permite contrastes complejos & Comparaciones no planificadas o contrastes complejos \\
\bottomrule
\end{tabular}
\caption{Comparación de métodos post-hoc}
\end{table}

\section{ANOVA de Dos Factores}

El ANOVA de dos factores (o de dos vías) permite estudiar simultáneamente el efecto de dos factores categóricos sobre una variable respuesta, así como su interacción.

\subsection{Modelo Estadístico}

El modelo para dos factores A (con $a$ niveles) y B (con $b$ niveles) es:

\begin{equation}
Y_{ijk} = \mu + \alpha_i + \beta_j + (\alpha\beta)_{ij} + \varepsilon_{ijk}
\end{equation}

donde:
\begin{itemize}
    \item $\mu$ es la media global
    \item $\alpha_i$ es el efecto principal del factor A (nivel $i$)
    \item $\beta_j$ es el efecto principal del factor B (nivel $j$)
    \item $(\alpha\beta)_{ij}$ es el efecto de interacción entre A y B
    \item $\varepsilon_{ijk} \sim N(0, \sigma^2)$ es el error aleatorio
    \item $k = 1, 2, \ldots, n_{ij}$ (réplicas en la celda $ij$)
\end{itemize}

\subsection{Efectos Principales e Interacción}

\begin{itemize}
    \item \textbf{Efecto principal del factor A:} Diferencia promedio entre niveles de A, ignorando B
    \item \textbf{Efecto principal del factor B:} Diferencia promedio entre niveles de B, ignorando A
    \item \textbf{Interacción A$\times$B:} El efecto de un factor depende del nivel del otro factor
\end{itemize}

\textbf{Interpretación de la interacción:}
\begin{itemize}
    \item Si la interacción es significativa, los efectos principales pueden no ser informativos por sí solos
    \item Gráficamente, las líneas de perfiles no son paralelas cuando hay interacción
    \item Si no hay interacción, las líneas son aproximadamente paralelas
\end{itemize}

\subsection{Tabla ANOVA de Dos Factores}

\begin{table}[H]
\centering
\small
\begin{tabular}{lccc}
\toprule
Fuente & Suma de Cuadrados & GL & Cuadrados Medios \\
\midrule
Factor A & $\text{SS}_A$ & $a-1$ & $\text{MS}_A = \text{SS}_A/(a-1)$ \\
Factor B & $\text{SS}_B$ & $b-1$ & $\text{MS}_B = \text{SS}_B/(b-1)$ \\
Interacción A$\times$B & $\text{SS}_{AB}$ & $(a-1)(b-1)$ & $\text{MS}_{AB} = \text{SS}_{AB}/[(a-1)(b-1)]$ \\
Error & $\text{SS}_E$ & $N-ab$ & $\text{MS}_E = \text{SS}_E/(N-ab)$ \\
Total & $\text{SS}_T$ & $N-1$ & \\
\bottomrule
\end{tabular}
\caption{Tabla ANOVA de dos factores}
\end{table}

Los estadísticos F para cada hipótesis son:
\begin{align*}
F_A &= \frac{\text{MS}_A}{\text{MS}_E} \sim F_{a-1, N-ab} \\
F_B &= \frac{\text{MS}_B}{\text{MS}_E} \sim F_{b-1, N-ab} \\
F_{AB} &= \frac{\text{MS}_{AB}}{\text{MS}_E} \sim F_{(a-1)(b-1), N-ab}
\end{align*}

\section{Diseños Experimentales Básicos}

Un diseño experimental bien planificado permite obtener conclusiones causales válidas y maximizar la eficiencia del estudio.

\subsection{Principios Básicos del Diseño Experimental}

\begin{enumerate}
    \item \textbf{Aleatorización:} Asignación aleatoria de tratamientos a unidades experimentales para evitar sesgos
    \item \textbf{Replicación:} Repetir el experimento para estimar el error experimental y aumentar la precisión
    \item \textbf{Control local (bloqueo):} Agrupar unidades experimentales homogéneas para reducir variabilidad
\end{enumerate}

\subsection{Diseño Completamente Aleatorizado (DCA)}

El DCA es el diseño más simple. Los tratamientos se asignan aleatoriamente a todas las unidades experimentales sin restricciones.

\textbf{Características:}
\begin{itemize}
    \item Máxima flexibilidad en la asignación
    \item Útil cuando las unidades experimentales son homogéneas
    \item Modelo: ANOVA de un factor
\end{itemize}

\textbf{Ventajas:} Análisis sencillo, máximos grados de libertad para el error. \\
\textbf{Desventajas:} Si hay heterogeneidad no controlada, aumenta el error experimental.

\subsection{Diseño en Bloques Completos Aleatorizados (DBCA)}

El DBCA agrupa las unidades experimentales en bloques homogéneos. Dentro de cada bloque, los tratamientos se asignan aleatoriamente.

\textbf{Características:}
\begin{itemize}
    \item Cada bloque contiene todos los tratamientos (una réplica completa)
    \item Se controla una fuente de variabilidad conocida (bloques)
    \item Modelo: ANOVA de dos factores sin interacción (bloques como factor)
\end{itemize}

\textbf{Modelo estadístico:}
\begin{equation}
Y_{ij} = \mu + \tau_i + \beta_j + \varepsilon_{ij}
\end{equation}

donde $\tau_i$ es el efecto del tratamiento $i$ y $\beta_j$ es el efecto del bloque $j$.

\textbf{Ventajas:} Reduce el error experimental al controlar variabilidad de bloques. \\
\textbf{Cuándo usar:} Cuando se puede identificar una fuente clara de variabilidad (ej: lotes de materia prima, días de medición, posición geográfica).

\subsection{Diseño en Cuadrado Latino}

El Cuadrado Latino controla dos fuentes de variabilidad simultáneamente mediante filas y columnas. Cada tratamiento aparece exactamente una vez en cada fila y una vez en cada columna.

\textbf{Características:}
\begin{itemize}
    \item Número de tratamientos = número de filas = número de columnas
    \item Controla dos fuentes de variabilidad
    \item Modelo: ANOVA de tres factores (tratamientos, filas, columnas) sin interacciones
\end{itemize}

\textbf{Ejemplo:} Estudio de rendimiento de 4 máquinas, controlando el efecto del operador (filas) y el turno (columnas).

\textbf{Ventajas:} Muy eficiente cuando hay dos fuentes de variabilidad. \\
\textbf{Desventajas:} Requiere que el número de niveles de todos los factores sea igual; no permite estimar interacciones.

\section{ANOVA de Medidas Repetidas}

El ANOVA de medidas repetidas se utiliza cuando se mide la misma unidad experimental en múltiples ocasiones (por ejemplo, a lo largo del tiempo).

\subsection{Modelo y Supuestos}

El modelo básico es:

\begin{equation}
Y_{ij} = \mu + \pi_i + \tau_j + \varepsilon_{ij}
\end{equation}

donde $\pi_i$ es el efecto del sujeto $i$ y $\tau_j$ es el efecto del tiempo (o condición) $j$.

\textbf{Supuesto de esfericidad:} Las varianzas de las diferencias entre todos los pares de mediciones repetidas deben ser iguales. Este supuesto es más restrictivo que la homocedasticidad del ANOVA clásico.

\subsection{Verificación de Esfericidad}

El test de Mauchly evalúa la hipótesis de esfericidad. Si se viola este supuesto, los grados de libertad del estadístico F deben corregirse:

\begin{itemize}
    \item \textbf{Corrección de Greenhouse-Geisser:} Más conservadora, recomendada cuando $\epsilon < 0.75$
    \item \textbf{Corrección de Huynh-Feldt:} Menos conservadora, recomendada cuando $\epsilon > 0.75$
\end{itemize}

donde $\epsilon$ es el estimador de esfericidad.

\section{Implementación Computacional}

\subsection{ANOVA en Python}

Python ofrece varias librerías para realizar ANOVA, principalmente \texttt{scipy} y \texttt{statsmodels}.

\subsubsection{ANOVA de un factor con scipy}

\begin{lstlisting}[style=python, caption=ANOVA de un factor en Python]
import numpy as np
from scipy import stats
import pandas as pd

# Datos de ejemplo: tiempo de reaccion (ms) para 3 tratamientos
grupo1 = [23, 25, 27, 29, 24]
grupo2 = [30, 32, 35, 33, 31]
grupo3 = [28, 26, 29, 30, 27]

# ANOVA de un factor
F_stat, p_value = stats.f_oneway(grupo1, grupo2, grupo3)

print(f"Estadistico F: {F_stat:.4f}")
print(f"p-valor: {p_value:.4f}")

if p_value < 0.05:
    print("Rechazamos H0: Al menos una media difiere")
else:
    print("No rechazamos H0: No hay diferencias significativas")
\end{lstlisting}

\subsubsection{ANOVA con statsmodels y post-hoc}

\begin{lstlisting}[style=python, caption=ANOVA completo con tabla y post-hoc]
import statsmodels.api as sm
from statsmodels.formula.api import ols
from statsmodels.stats.multicomp import pairwise_tukeyhsd

# Crear DataFrame
data = pd.DataFrame({
    'tiempo': grupo1 + grupo2 + grupo3,
    'tratamiento': ['A']*5 + ['B']*5 + ['C']*5
})

# Ajustar modelo ANOVA
modelo = ols('tiempo ~ C(tratamiento)', data=data).fit()
tabla_anova = sm.stats.anova_lm(modelo, typ=2)

print(tabla_anova)

# Calcular eta cuadrado
ssb = tabla_anova.loc['C(tratamiento)', 'sum_sq']
sst = tabla_anova['sum_sq'].sum()
eta_squared = ssb / sst
print(f"\nEta cuadrado: {eta_squared:.4f}")

# Test post-hoc de Tukey
tukey = pairwise_tukeyhsd(endog=data['tiempo'], 
                          groups=data['tratamiento'], 
                          alpha=0.05)
print("\nTest de Tukey HSD:")
print(tukey)
\end{lstlisting}

\subsubsection{ANOVA de dos factores}

\begin{lstlisting}[style=python, caption=ANOVA de dos factores con interacción]
# Datos con dos factores
data2 = pd.DataFrame({
    'rendimiento': [45, 47, 50, 52, 43, 46, 55, 58, 60, 62, 54, 57],
    'fertilizante': ['A', 'A', 'A', 'A', 'B', 'B', 'B', 'B', 
                     'C', 'C', 'C', 'C'],
    'riego': ['bajo', 'bajo', 'alto', 'alto', 'bajo', 'bajo', 
              'alto', 'alto', 'bajo', 'bajo', 'alto', 'alto']
})

# ANOVA de dos factores con interacción
modelo2 = ols('rendimiento ~ C(fertilizante) * C(riego)', 
              data=data2).fit()
tabla_anova2 = sm.stats.anova_lm(modelo2, typ=2)

print(tabla_anova2)

# Grafico de interacción
import matplotlib.pyplot as plt
import seaborn as sns

plt.figure(figsize=(8, 5))
sns.pointplot(data=data2, x='fertilizante', y='rendimiento', 
              hue='riego', markers=['o', 's'], capsize=0.1)
plt.title('Grafico de Interaccion')
plt.ylabel('Rendimiento')
plt.xlabel('Fertilizante')
plt.legend(title='Riego')
plt.grid(True, alpha=0.3)
plt.tight_layout()
plt.savefig('interaccion.png', dpi=150)
plt.show()
\end{lstlisting}

\subsection{ANOVA en R}

R tiene funciones nativas muy potentes para ANOVA, incluyendo \texttt{aov()}, \texttt{lm()}, y el paquete \texttt{car}.

\subsubsection{ANOVA de un factor en R}

\begin{lstlisting}[style=rstyle, caption=ANOVA de un factor en R]
# Crear datos
grupo1 <- c(23, 25, 27, 29, 24)
grupo2 <- c(30, 32, 35, 33, 31)
grupo3 <- c(28, 26, 29, 30, 27)

# Crear data frame
datos <- data.frame(
  tiempo = c(grupo1, grupo2, grupo3),
  tratamiento = factor(rep(c("A", "B", "C"), each = 5))
)

# ANOVA
modelo <- aov(tiempo ~ tratamiento, data = datos)
summary(modelo)

# Verificar supuestos
# Normalidad de residuos
shapiro.test(residuals(modelo))

# Homocedasticidad
library(car)
leveneTest(tiempo ~ tratamiento, data = datos)

# Graficos de diagnostico
par(mfrow = c(2, 2))
plot(modelo)

# Test post-hoc de Tukey
TukeyHSD(modelo)
\end{lstlisting}

\subsubsection{ANOVA de dos factores en R}

\begin{lstlisting}[style=rstyle, caption=ANOVA de dos factores en R]
# Datos
datos2 <- data.frame(
  rendimiento = c(45, 47, 50, 52, 43, 46, 55, 58, 60, 62, 54, 57),
  fertilizante = factor(rep(c("A", "B", "C"), each = 4)),
  riego = factor(rep(c("bajo", "alto"), 6))
)

# ANOVA de dos factores con interacción
modelo2 <- aov(rendimiento ~ fertilizante * riego, data = datos2)
summary(modelo2)

# Grafico de interacción
interaction.plot(x.factor = datos2$fertilizante,
                 trace.factor = datos2$riego,
                 response = datos2$rendimiento,
                 fun = mean,
                 type = "b",
                 col = c("blue", "red"),
                 pch = c(19, 17),
                 lty = 1,
                 lwd = 2,
                 xlab = "Fertilizante",
                 ylab = "Rendimiento medio",
                 trace.label = "Riego")
\end{lstlisting}

\subsubsection{Diseño en Bloques Completos Aleatorizados}

\begin{lstlisting}[style=rstyle, caption=DBCA en R]
# Datos de DBCA
datos_dbca <- data.frame(
  produccion = c(50, 55, 58, 48, 53, 60, 52, 57, 62),
  maquina = factor(rep(c("M1", "M2", "M3"), 3)),
  operador = factor(rep(c("Op1", "Op2", "Op3"), each = 3))
)

# ANOVA para DBCA (bloques como factor, sin interacción)
modelo_dbca <- aov(produccion ~ maquina + operador, data = datos_dbca)
summary(modelo_dbca)

# Comparaciones múltiples para máquinas
TukeyHSD(modelo_dbca, "maquina")
\end{lstlisting}

\section{Ejercicios Propuestos}

\subsection{Ejercicio 1: ANOVA de un factor}

Un investigador desea comparar la efectividad de cuatro métodos de enseñanza (A, B, C, D) midiendo el puntaje en un examen estandarizado. Los datos son:

\begin{itemize}
    \item Método A: 78, 82, 79, 85, 80
    \item Método B: 88, 90, 85, 92, 87
    \item Método C: 75, 78, 73, 80, 76
    \item Método D: 82, 84, 86, 83, 85
\end{itemize}

\textbf{Tareas:}
\begin{enumerate}
    \item Realizar un ANOVA de un factor ($\alpha = 0.05$)
    \item Verificar los supuestos (normalidad, homocedasticidad)
    \item Calcular $\eta^2$ e interpretar
    \item Si hay diferencias significativas, realizar un test post-hoc de Tukey
\end{enumerate}

\subsection{Ejercicio 2: ANOVA de dos factores}

Un estudio evalúa el efecto de dos factores en el tiempo de secado de pintura:
\begin{itemize}
    \item Factor A (tipo de pintura): látex, acrílica
    \item Factor B (temperatura): baja (15°C), media (25°C), alta (35°C)
\end{itemize}

Se miden tiempos de secado (minutos) con 3 réplicas por combinación de tratamientos.

\textbf{Tareas:}
\begin{enumerate}
    \item Realizar ANOVA de dos factores
    \item Evaluar efectos principales e interacción
    \item Graficar el diagrama de interacción e interpretar
    \item Si hay interacción significativa, analizar efectos simples
\end{enumerate}

\subsection{Ejercicio 3: Diseño en Bloques}

Se estudia el rendimiento de 4 variedades de trigo en 5 parcelas (bloques). Cada bloque contiene las 4 variedades asignadas aleatoriamente.

\textbf{Tareas:}
\begin{enumerate}
    \item Identificar el diseño experimental
    \item Realizar el ANOVA correspondiente
    \item Comparar con un DCA (ignorando bloques) y comentar diferencias
    \item Realizar comparaciones múltiples entre variedades
\end{enumerate}

\subsection{Ejercicio 4: Análisis Completo}

Datos simulados de un experimento sobre productividad laboral bajo diferentes niveles de ruido (silencio, ruido bajo, ruido moderado, ruido alto).

\textbf{Tareas:}
\begin{enumerate}
    \item Análisis exploratorio: boxplots, estadísticas descriptivas
    \item Verificar supuestos de ANOVA
    \item Realizar ANOVA e interpretar resultados
    \item Calcular tamaño del efecto
    \item Realizar comparaciones post-hoc apropiadas
    \item Redactar conclusiones en lenguaje no técnico
\end{enumerate}

\section{Conclusiones}

El ANOVA es una herramienta fundamental en estadística aplicada que permite comparar medias de múltiples grupos de manera rigurosa. Los conceptos clave incluyen:

\begin{itemize}
    \item La descomposición de variabilidad en componentes explicados y no explicados
    \item La importancia de verificar supuestos para validez de las conclusiones
    \item El uso de comparaciones múltiples post-hoc para identificar diferencias específicas
    \item La extensión a diseños más complejos (dos factores, medidas repetidas)
    \item El diseño experimental adecuado como base para inferencias causales válidas
\end{itemize}

En la práctica, el ANOVA se complementa con visualizaciones efectivas, verificación cuidadosa de supuestos, y reporte transparente de todos los análisis realizados. Las implementaciones en Python y R facilitan estos análisis, permitiendo al analista enfocarse en la interpretación sustantiva de los resultados.

\vspace{1cm}

\noindent\textbf{Referencias recomendadas:}
\begin{itemize}
    \item Montgomery, D. C. (2017). \textit{Design and Analysis of Experiments}. Wiley.
    \item Kutner, M. H., et al. (2005). \textit{Applied Linear Statistical Models}. McGraw-Hill.
    \item Field, A., Miles, J., \& Field, Z. (2012). \textit{Discovering Statistics Using R}. SAGE.
\end{itemize}

\end{document}
